
\section{Issue Tracking}
Si definisce \textbf{Atlassian Jira} come principale strumento di issue tracking del team.
\textbf{Jira} diventa inoltre unica fonte di verità sullo stato del progetto e delle attività in corso, il cui contenuto
viene riassunto ed esportato nei verbali e diari di bordo del team.

Come visibile nelle procedure successivamente definite, \textit{Jira} è al centro del lavoro del team.

\subsection{Progetti (\textit{spaces}) in Jira}

Per ogni scope oggetto dell'attenzione del team, si crea uno \textit{space} apposito in Jira.
Lo space apposto all'amministrazione del team e la documentazione conseguente è identificato dal nome (e chiave) "\textit{SWE}".
Per lo space \textit{SWE}, si sceglie la modalità "Kanban" di facile gestione e adatta alla natura dello space.

A partire dall'aggiudicazione del capitolato, viene creato uno spazio apposito al progetto in Jira.
Lo spazio sarà utilizzato per la definizione di qualsiasi attività interna al progetto.